\section{Lebensohl-ish, and good-bad 2NT}

\subsection{Lebensohl}

It is widely agreed that distinguishing ranges using 2N as artificial (i.e. Lebensohl) is very beneficial.
However, we need to define the "cases" carefully:

\begin{itemize}
    \setlength\itemsep{0pt}
    \item this only applies to advancer or responder
    \item forced 2N by t/o = normal Lebensohl
    \item (P/bid) - 1N - (2X) - 2N or (1Y) - X - (2X/2Y) - 2N = transfer Lebensohl. (because these are free bids)
    \item t/o twice, or resp shows 0-5, or t/o'er passed, 2N = scramble
    \item other free 2N = nat unless otherwise specified
\end{itemize}

\bin{\currfiledir/seq/Leb}

\subsection{good-bad 2NT}

Similar for opener's rebid, in competition, using 2N to distinguish ranges are also common practices. 
We also need to define the cases carefully: suppose the auction goes \texttt{1X - (A) - B - (C)}

\begin{itemize}
    \setlength\itemsep{0pt}
    \item B should never be P. i.e. if resp passes, there are never good-bad. (note: if resp passes, then 2N is unusual)
    \item (1) C > 2X
    \item (2) 1X - (jump overcall) - X - (P)
    \item (3) 1S - (2H) - X - (P)
\end{itemize}

\bin{\currfiledir/seq/good_bad}

\subsection{extended good-bad 2NT}

\textit{not now, but maybe in the future}
